\documentclass[10pt,a4paper]{amsart}
\usepackage[margin=19mm]{geometry}
\usepackage{amsmath,amssymb,mathtools}
\usepackage[scr=rsfs,cal=euler]{mathalfa}
\usepackage{xfrac}
\usepackage[unicode,hidelinks,bookmarks=false]{hyperref}
\newcommand{\ie}{i.e.\ }
\newcommand{\cf}{cf.\ }
\newcommand{\resp}{respectively\ }
\newcommand{\Subsection}{\subsection}
\providecommand{\nobreakdash}{\nobreak\hbox{-}}
\setlength{\emergencystretch}{2em}
\multlinegap0pt
\usepackage{bm}
\usepackage{bbm}
\usepackage{dsfont}
\usepackage{xspace}
\usepackage{cancel}
\usepackage{mathtools}
\usepackage{amsthm}
\makeatletter
\@ifundefined{theorem}{\newtheorem{theorem}{Theorem}}{}
\makeatother
\renewcommand\S{\mathcal S}
\title[Arrow pattern avoidance in permutations: structure and enume]{Arrow pattern avoidance in permutations: structure and enumeration}
\author{Kassie Archer, Robert P. Laudone}
\date{Source: arXiv:2603.04218v1 (CC BY 4.0)}
\begin{document}
\maketitle

\noindent\emph{Source and licence.} Arrow pattern avoidance in permutations: structure and enumeration. Version arXiv:2603.04218v1 (CC BY 4.0), distributed under the Creative Commons Attribution 4.0 International licence, as recorded in the arXiv licence metadata for this exact version and shown on the abstract page https://arxiv.org/abs/2603.04218v1. Full rights notice: licenses/arxiv-2603.04218-CC-BY-4.0.txt.\par\smallskip

\noindent\textit{Editorial scope.} Counted unit: the Theorem whose source label is \texttt{thm:12t23} in the article (source0.tex lines 381--383, source label \texttt{thm:12t23}), delivered together with the article's own complete original proof of it (source0.tex lines 385--416, one \texttt{proof} environment of 32 lines). No mathematics is written, repaired or re-derived: statement and proof are the article's own text line for line. Required context retained uncounted: none. Every definition, display and label of those lines is retained unchanged. Omitted from this package and not redistributed: 1--380, 384--384, 417--854. Nothing inside a retained excerpt is omitted or altered: every retained body is delivered exactly as the article's lines are written, so no omission receipt of any kind accompanies this package. Citations: the retained excerpts carry no citation command at all, so no bibliography entry of the article is transcribed and no reference is redistributed. Reference adaptation: none was needed and none was made; every reference inside the delivered unit resolves inside it, so no unresolved reference or citation is produced. Printed slips kept literal: no printed word, symbol, hypothesis or number of the counted statement or of its proof was altered. Layout and class: the article's own class is \texttt{article} with packages including authblk, hyperref, float, amssymb, amsthm, amsmath, mathrsfs, fullpage, multirow, tikz, enumerate; the wrapper is the lane's pinned \texttt{amsart} editorial wrapper at 10pt on A4, which restates the theorem-like environments the retained text needs and adds the article's own macro definitions that the retained text needs (\texttt{\textbackslash S}), with extra wrapper packages (bm, bbm, dsfont, xspace, cancel, mathtools). The delivered numbering is the unit's own. Assets: no retained span contains a figure environment, a graphics-inclusion command, a PDF-page inclusion, a table or a TikZ picture; no image file is copied into this package, the package declares no asset dependency and no image file is redistributed with it. Licence: version arXiv:2603.04218v1 of this article is distributed under the Creative Commons Attribution 4.0 International licence, as recorded in the arXiv licence metadata for this exact version and shown on the abstract page https://arxiv.org/abs/2603.04218v1. A full rights notice is delivered at licenses/arxiv-2603.04218-CC-BY-4.0.txt. Characters: every retained excerpt is pure ASCII, so the delivered engine, XeLaTeX with its default Latin Modern text font, reports no missing character.
\begin{theorem} \label{thm:12t23}
    For $n \geq 2$, $a_n(12; 2 \to 3) = 2B_{n-1}$ and $a_1(12; 2 \to 3) = 1$.
\end{theorem}

\begin{proof}
    Suppose $n \geq 3$ and $\pi \in \S_n(12; 2 \to 3)$. Let $b_n$ denote the number of these permutations with $\pi_n = 1$. First, notice that $a_n = a_{n-1} + b_n$. Indeed, $\pi \in \S_n(12; 2 \to 3)$ either has $\pi_n = n$ or $\pi_n = 1$. If not, then $\pi_n = k$ for some $2 \leq k \leq n-1$. Suppose $\pi_j = 1$ and $\pi_\ell = n$, then $(\pi_j\pi_n;\pi_n \to \pi_\ell)$ is a $(12; 2 \to 3)$ pattern. There are $a_{n-1}$ such permutations with $\pi_n = n$ and by definition $b_n$ with $\pi_n = 1$.

    We will now find a recursion for $b_n$. So for now, we only consider $\pi \in \S_n(12;2\to3)$ with $\pi_n = 1$. We break these permutations up into two groups. The first ends with $n > \pi_{i_1} > \pi_{i_2}> \ldots > \pi_{i_k} > 1$. The second group has some ascent between $n$ and $1$. 

    To enumerate the first group, we can select any $k$ elements from $2,\dots,n-1$ to be the decreasing portion between $n$ and $1$. The elements before $n$ can then be arranged in $a_{n-k-2}$ ways. So the first group of permutations is enumerated by $\sum_{k=0}^{n-2} \binom{n-2}{k} a_{n-k-2}$.

    To enumerate the second group, consider the final ascent that occurs between $n$ and $1$. Say this ascent is $i < j$ so that $\pi = \pi_1 \pi_2 \cdots \pi_m ij \pi_{m+3} \cdots \pi_n$ with $\pi_{m+3} > \cdots > \pi_n = 1.$ To avoid $(12;2\to3)$ we must have all the elements less than $i$ occurring to the right of $j$. If not, there is some element $k < i$ to the left of $j$, and therefore $i$, so $(ki; i \to j)$ is a $(12; 2\to 3)$ pattern. Notice this only works because we necessarily have $i \to j$ since the ascent appears after $n$. But since $i < j$ is the final ascent in $\pi$, all elements after $j$ are decreasing, so $\pi$ must end with $(i-1)(i-2)\ldots21$.
    
    If we fix $2 \leq i \leq n-2$ that begins the ascent, we can enumerate these permutations by choosing the elements between $i$ and $i-1$ in $\pi$, which must be decreasing, non-empty and in $\{i+1,\dots,n-1\}$. If there are $k$ such elements, we can choose them in $\binom{n-i-1}{k}$. The elements to the left of $i$ are enumerated by $b_{n-i-k+1}$ since ${\rm red}(\pi_1 \cdots \pi_m i)$ must be in $\S_n(12; 2\to3)$ and  $i$ reduces to $1$.
    
    For example, the permutation $784{\bf 65} 321 \in \S_7(12; 2\to3)$ has final ascent beginning with $i = 4$, we then choose $6$ and $5$ to appear after $4$, and are forced to end the permutation with $321$. After removing all the elements to the right of $i$ and reducing we find ${\rm red}(784) = 231 \in \S_3(12; 2\to3)$ with final entry $1$.

    Summing over all choices of $i$, there are
    \[
    \sum_{i=2}^{n-2} \sum_{k=1}^{n-i-1} \binom{n-i-1}{k} b_{n-i-k+1}
    \]
    such permutations. This means
    \[
    b_n = \sum_{k=0}^{n-2} \binom{n-2}{k} a_{n-k-2} + \sum_{i=2}^{n-2} \sum_{k=1}^{n-i-1} \binom{n-i-1}{k} b_{n-i-k+1}.
    \]
    We can simplify the first and second sums to find:
    \[
    b_n = -a_{n-2} + a_1 + \sum_{j=0}^{n-2} \binom{n-2}{j} a_{j} + \sum_{m=2}^{n-2} \left[\binom{n-2}{m-1}\right]b_m.
    \]
    Since $b_n = a_n - a_{n-1}$, we can express this as a recursion in $a_n$,
    \begin{align*}
    a_n &= a_{n-1} - a_{n-2} + a_1 + \sum_{j=0}^{n-2} \binom{n-2}{j} a_j + \sum_{m=2}^{n-2} \binom{n-2}{m-1}(a_m - a_{m-1})\\
    &= 1 + \sum_{m=1}^{n-1} \binom{n-2}{m-1} a_m
    \end{align*}
    This recursion implies the desired result.
\end{proof}

\end{document}
